\documentclass[11pt]{article}
\usepackage[left=1in, right=1in, bottom=1in, top=1.5in]{geometry}
\usepackage{amsmath}
%\usepackage[mathscr]{euscript} %for \mathscr{}
\usepackage{amssymb}  %for \mathbb{}
\usepackage{mathrsfs} %for \mathscr{}
\usepackage{hyperref} %for \href{}{}
\usepackage{multirow} %for \multirow{2}{*}{hola}
\usepackage{amsthm}   %for \theoremstyle, \begin{proof}
\usepackage{cite}     %for multiple cite as \cite{cite1, cite2}
 \usepackage{relsize} %for \mathlarger

\theoremstyle{definition}     % for put not italic letters
\newtheorem{theorem}{Theorem} %[section]         % Teoremas
\theoremstyle{definition}     % for put not italic letters
\newtheorem{lemma}{Lemma}     %[section]         % Teoremas
\hypersetup{colorlinks=true,linkcolor=blue,urlcolor=cyan}
\def\bigsum{\mathlarger{\mathlarger{\sum}}}
\def\bigmid{\,\mathlarger{\mathlarger{\mathlarger{\mathlarger{\mid}}}}\,}
\def\be{\begin{equation}}
\def\ee{\end{equation}}
\def\Re{\mathbb{R}}
\newcommand{\partd}[2]{\frac{\partial #1}{\partial #2}}

\begin{document}

\title{Relativistic Statistical Mechanics}
\author{Juan A. Guarín-Rojas \\
María Fernanda Carvajal \\
Carlos José Páez González \\
Antonio Calixto Gutiérrez Piñerez }
\date{\today}
\maketitle

%\begin{center}
%	\textsc{Abstract} Hola
%\end{center}

\section{Transformation Laws Table}
The table \ref{Tabla1} shows the results of the literature \cite{einstein, planck, Ott, arzelies1965transformation, Sutcliffe} and the results of this work for the transformation laws of the thermodynamic variables. Our results agree with \cite{Sutcliffe} and partially agree with \cite{Ott, arzelies1965transformation}. 

\begin{table}[h!]
   \centering
    \begin{tabular}{l|c|c|c|c}
        \hline
                    & \multirow{2}{7em}{\textbf{Einstein, Planck}.  \cite{einstein, planck} } & \multirow{2}{7em}{\textbf{Ott, Arzelièz}. \cite{Ott, arzelies1965transformation} } & \multirow{2}{*}{\textbf{Sutcliffe} \cite{Sutcliffe}}  & \multirow{2}{9em}{\textbf{Guarín-Carvajal-Páez-Gutiérrez}} \\ 
                    & & & & \\ \hline
        \multirow{2}{4em}{Pressure} & \multicolumn{2}{c|}{\multirow{2}{*}{$P=P_0$}} & \multirow{2}{5em}{$P_m=P_0$, $P_{th}=P_0\gamma^2$} & \multirow{2}{*}{$P=P_0\gamma^2$} \\	
        & \multicolumn{2}{c|}{ } & & \\ \cline{1-5}
        Volume      & \multicolumn{4}{c}{$V=V_0/\gamma$}                                  \\ 
        Entropy     & \multicolumn{4}{c}{$S=S_0$}                                         \\ \cline{1-5}

        Temperature & $T=T_0/\gamma$                  & \multicolumn{3}{c}{$T=T_0\gamma$}\\ %\cline{2-5}
        Heat            & $\delta Q =  \delta Q_0/\gamma$  & \multicolumn{3}{c}{$\delta Q = \delta Q_0\gamma$} \\ %\cline{2-5}
        Internal energy & $E=E_0\gamma+P_0V_0\gamma v^2/c^2$ & \multicolumn{3}{c}{$E=E_0\gamma$} \\ \hline
        \multirow{2}{5em}{Helmholtz free energy} & \multirow{2}{*}{$F =   F_0/\gamma$} & \multicolumn{3}{c}{\multirow{2}{*}{$F=F_0\gamma$}} \\
        & \\ \hline
    \end{tabular}
    \label{Tabla1}
    \caption{Transformation of the thermodynamic quantities for different authors. This table was adapted from \cite{Balescu}. A zero subindex indicates variables measured for a rest observer with the container of the gas, by convention $\gamma=1/\sqrt{1-v^2/c^2}$ is the Lorentz factor, $P_m$ is the \textit{mechanical pressure}, and $P_{th}$ the \textit{thermodynamic} one (as discussed by \cite{Sutcliffe}).
    }
\end{table}

\section{The invariance of entropy}
\begin{theorem}
The entropy $S=k_B\ln\Omega$ for a moving inertial observer $K$ with respect to the container of an ideal gas ($N$ equal particles of same mass, elastic collisions between particles and walls, and non-interacting particles) is equal to the entropy $S_0=k_B\ln\Omega_0$ for a rest observer $K_0$ with the container. 

Here, $\Omega=\Omega(E,V,N)$ is the number of states for $K$ depending on the internal energy, the volume, and the number of particles respectively. Analogously, $\Omega_0=\Omega_0(E_0,V_0,N_0)$ is the number of states for $K_0$ depending on its corresponding internal energy, volume and number of particles respectively. Besides, the number of particles for each frame must be the same, this is $N=N_0$.
\end{theorem}

\textit{Proof:}

By definition
\be 
\Omega=\frac{1}{N!}\frac{1}{h^{3N}}\underbrace{\int\limits_\mathscr{V} d^3r_1\cdots d^3r_N}_{\Omega_v}\,\underbrace{\int\limits_\mathcal{R} d^3p_1 \cdots d^3 p_N}_{\Omega_p}\,,
\ee
where $h$ is the planck constant, $\mathbf{r}_i$ and $\mathbf{p}_i$ are the $3-$position and $3-$momentum of the ith particle for $K$, $\mathscr{V}$ is the volume space, and $\mathcal{R}$ is the momentum space defined as
\be 
\mathcal{R} = \left\{ (\mathbf{g}_1,\dots,\mathbf{g}_N)\in\Re^{3N} \mid \sum\limits_{i=1}^{N} c\sqrt{m^2c^2+|\mathbf{g}_i|^2} = E_0 \right\} \,, \label{eq:R}
\ee
where $\mathbf{g}_i$ is the $3-$momentum of the ith particle according to $K_0$. Now, by the non-interacting condition, the position $\mathbf{r}_i$ of each particle is independent one of the other. And therefore we can separate $\Omega_v$ as a product of integrals
\be 
	\Omega_v = \left(\int\limits_V d^3 r_1\right)\cdots \left(\int\limits_V d^3 r_N\right) = \left(\int\limits_V d^3 r_1\right)^N = V^N\,.
\ee

\begin{lemma}
The volume $V_0$ for the observer $K_0$ changes to $V_0/\gamma=V$ by a Lorentz transformation corresponding to a Lorentz Boost in one direction. $\gamma$ is the Lorentz factor.
\end{lemma}

By the previous lemma 
\be 
\Omega_v = \frac{V_0\,^N}{\gamma^N} = \frac{\Omega_{v_0}}{\gamma^N}\,. \label{eq:omega_v}
\ee

By the other side, chosen an appropriate coordinate axis we can say that the Lorentz Boost occur in the $x-$axis. According to \cite{Landau} the transformation laws for momentum in this case are 
\begin{eqnarray}
	p_{x_i} &=& \gamma\left(g_{x_i} + \frac{v}{c^2}E_{0_i}\right) \,, \\
	p_{y_i} &=& g_{y_i} \,, \\
	p_{z_i} &=& g_{z_i} \,,
\end{eqnarray}
where $\mathbf{p}_i$ is the momentum of the ith particle for observer $K$, and $\mathbf{g}_i$ is the momentum of the same particle but for $K_0$. Besides, $v$ is the relative speed between the observers, $\gamma=1/\sqrt{1-v^2/c^2}$, $c$ is the speed of light, and $E_{0_i}=c\sqrt{m^2c^2+g_i\,^2}=mc^2\sqrt{1+\frac{g_i\,^2}{m^2c^2}}$ is the relativistic energy of the particle for $K_0$.

We want to implement a change of variables from $p_i$ to $g_i$ in the momentum integrals. The Jacobian of the transformation is 
\be 
	J(\mathbf{g}) = \frac{\partial p_{x^i}}{\partial g_{x^i}} =  \begin{vmatrix}
     \partd{p_x}{g_x} & \partd{p_x}{g_y} & \partd{p_x}{g_z}\\ 
     0 & 1 & 0\\
     0 & 0 & 1  
\end{vmatrix} = \partd{p_x}{g_x} = \partd{ }{g_x}\left[ \gamma \left(g_x + \frac{v}{c^2}mc^2\sqrt{1+\frac{g^2}{m^2c^2}} \right)\right]\,.
\ee
Knowing that $g^2=g_x\,^2+g_y\,^2+g_z\,^2$,
\be 
	J(\mathbf{g}) = \gamma \left( 1 + \frac{v}{c^2}mc^2 \frac{g_x/(m^2c^2)}{\sqrt{1+\frac{g^2}{m^2c^2}}} \right) = \gamma \left( 1 + \frac{v}{c} \frac{g_x/(mc)}{\sqrt{1+\frac{g^2}{m^2c^2}}} \right)\,.
\ee
Then by implement this coordinate transformation for all the variables of the $N$ particles one gets
\be 
	\Omega_p = \int\limits_\mathcal{R} d^3p_1 \cdots d^3p_N = \int\limits_\mathcal{R} \gamma^N \prod\limits_{i=1}^{N}\left( 1 + \frac{v}{c} \frac{g_{x_i}/(mc)}{\sqrt{1+\frac{g_i\,^2}{m^2c^2}}} \right) d^3 g_1 \cdots d^3g_N\,. \label{eq:omega_p}
\ee
On the other side, one can check that
\begin{eqnarray}
	\prod\limits_{i=1}^{N}(1+M_i) &=& 1+M_1+M_2+\cdots+M_N+ \notag  \\  
								 & & +M_1 M_2 +M_1M_3+\cdots+M_1M_N+\cdots M_{N-1}M_n+ \notag \\
								 & & +M_1M_2M_3 +\cdots M_1 M_{N-1}M_N + \cdots M_{N-2}M_{N-1}M_N + \cdots\,, \notag \\
								 &=& 1 + \sum\limits_{k=1}^{N} \bigsum \limits_{1\leq i_1 < i_2 < \cdots < i_k\leq N} M_{i_1}\cdots M_{i_k}\,,
\end{eqnarray}
where $\sum_{1\leq i_1 < i_2 < \cdots < i_k\leq N}$ stands for the sum over all integers $i_1,\dots,i_k$ that satisfies $1\leq i_1 < i_2 < \cdots < i_k\leq N$. Therefore
\be 
	\Omega_p = \int\limits_\mathcal{R} \gamma^N \left( 1 + \sum\limits_{k=1}^{N} \bigsum\limits_{1\leq i_1 < i_2 < \cdots < i_k\leq N} \frac{v^k}{c^k} \frac{g_{x_{i_1}}/(mc)}{\sqrt{1+\frac{g_{i_i}\,^2}{m^2c^2}}} \cdots \frac{g_{x_{i_k}}/(mc)}{\sqrt{1+\frac{g_{i_k}\,^2}{m^2c^2}}} \right) d^3 g_1 \cdots d^3g_N\,.
\ee
Now, as integration is a linear operation,
\be 
	\Omega_p =  \gamma^N \int\limits_\mathcal{R} d^3 g_1 \cdots d^3g_N  + \gamma^N \sum\limits_{k=1}^{N} \bigsum\limits_{1\leq i_1 < i_2 < \cdots < i_k\leq N} \frac{v^k}{c^k} \int\limits_\mathcal{R} \frac{g_{x_{i_1}}/(mc)}{\sqrt{1+\frac{g_{i_i}\,^2}{m^2c^2}}} \cdots \frac{g_{x_{i_k}}/(mc)}{\sqrt{1+\frac{g_{i_k}\,^2}{m^2c^2}}} d^3 g_1 \cdots d^3g_N\,. \label{eq:sums}
\ee

While integrating the second term we will find that if a point $(\mathbf{g}_1,\dots,\mathbf{g}_k)$ is in $\mathcal{R}$, then $(-\mathbf{g}_1,\mathbf{g}_2,\dots,\mathbf{g}_k)$ also will be in $\mathcal{R}$. This can be seen by looking in the definition of the momentum space \ref{eq:R}. Therefore, $\mathcal{R}$ is composed by pairs of points $(\mathbf{g}_1,\dots,\mathbf{g}_N)$ and $(-\mathbf{g}_1,\mathbf{g}_2,\dots,\mathbf{g}_N)$ such that its respective contribution to the integral cancel out each other, this is 
\be 
	\int\limits_\mathcal{R} \frac{g_{x_{i_1}}/(mc)}{\sqrt{1+\frac{g_{i_i}\,^2}{m^2c^2}}} \cdots \frac{g_{x_{i_k}}/(mc)}{\sqrt{1+\frac{g_{i_k}\,^2}{m^2c^2}}} d^3 g_1 \cdots d^3g_N = 0\,.
\ee
Then only the first term in \ref{eq:sums} remains
\be 
\	\Omega_p =  \gamma^N \int\limits_\mathcal{R} d^3 g_1 \cdots d^3g_N\,.
\ee

Regardless the factor $\gamma^N$, this integral is exactly the definition of $\Omega_p$ but for $K_0$. Hence $\Omega_p=\gamma^N \Omega_{p_0}$, and so $\Omega=\Omega_0$. This ends the proof.

%\section{Appendix}
%In local coordinates the components of the $4-$velocity $\vec{u}_i$ of a particle are $(\gamma(v_i),\gamma(v_i)\,\mathbf{v}_i/c)$, where $\mathbf{v}_i$ is the $3-$velocity of the particle seen by the local frame and $\gamma(v_i)=1/\sqrt{1-v_i\,^2/c^2}$. Besides, the components of its $4-$momentum $\vec{p}_i=m\vec{u}_i$ are $(E_i/c^2, \mathbf{p}_i/c)$, where $E_i$ is the relativistic energy of the particle seen by the local frame and $\mathbf{p}_i$ is the corresponding $3-$momentum seen by the local frame.

\bibliographystyle{unsrt}
\bibliography{BibTex}

\end{document}